\section{Details on Laurent power series}

Fix a commutative ring $K$.

\subsection{Proof of Theorem~\ref{thm.fps.laure.laupol-ring} (sketched)}

The set $K\left[\xpm\right]$ is closed under addition, scaling, and
the convolution product from Definition~\ref{def.fps.laure.laupol}
(analogous to Theorem~\ref{thm.fps.pol.ring}, replacing
$\sum_{i=0}^{n}$ sums by $\sum_{i \in \mathbb{Z}}$ sums).
The multiplication is associative, commutative, distributive, and
$K$-bilinear (analogous to the corresponding parts of
Theorem~\ref{thm.fps.ring}). The element
$\left(\delta_{i,0}\right)_{i \in \mathbb{Z}}$ is a neutral element
for multiplication.

It remains to show that $x$ is invertible. Set
$\overline{x} := \left(\delta_{i,-1}\right)_{i \in \mathbb{Z}}$.
Then $x \overline{x} = 1$ and $\overline{x} x = 1$ by direct
calculation. For example, from
$x = \left(\delta_{i,1}\right)_{i \in \mathbb{Z}}$ and
$\overline{x} = \left(\delta_{i,-1}\right)_{i \in \mathbb{Z}}$, we
get
\[
  x \overline{x}
  = \left(c_n\right)_{n \in \mathbb{Z}},
  \quad \text{where} \quad
  c_n = \sum_{i \in \mathbb{Z}} \delta_{i,1}\,\delta_{n-i,-1}
\]
(by Definition~\ref{def.fps.laure.laupol}).
For each $n \in \mathbb{Z}$,
\begin{align*}
  c_n
  &= \sum_{i \in \mathbb{Z}} \delta_{i,1}\,\delta_{n-i,-1}
   = \underbrace{\delta_{1,1}}_{= 1} \delta_{n-1,-1}
     + \sum_{\substack{i \in \mathbb{Z} \\ i \neq 1}}
       \underbrace{\delta_{i,1}}_{= 0\text{ (since } i \neq 1\text{)}}
       \delta_{n-i,-1} \\
  &= \delta_{n-1,-1} = \delta_{n,0},
\end{align*}
since $n - 1 = -1$ if and only if $n = 0$. Hence
$x \overline{x} = \left(\delta_{n,0}\right)_{n \in \mathbb{Z}} = 1$.
The proof of $\overline{x} x = 1$ is analogous, or follows by commutativity.

\subsection{Helpers for Theorem~\ref{thm.fps.laure.laupol-ring}}

\begin{lemma}[$T(1) \cdot T(-1) = 1$]
\label{lem.fps.laure.T-inv}
\lean{AlgebraicCombinatorics.FPS.Laurent.T_inv}
\leanhelper
\leanok
In the Laurent polynomial ring, $T(1) \cdot T(-1) = 1$.
This witnesses the right inverse of $x$.
\end{lemma}

\begin{proof}
\leanok
Follows from $T(a) \cdot T(b) = T(a+b)$ with $a = 1$, $b = -1$.
\end{proof}

\begin{lemma}[$T(-1) \cdot T(1) = 1$]
\label{lem.fps.laure.T-neg-one-mul-T}
\lean{AlgebraicCombinatorics.FPS.Laurent.T_neg_one_mul_T}
\leanhelper
\leanok
In the Laurent polynomial ring, $T(-1) \cdot T(1) = 1$.
This witnesses the left inverse of $x$.
\end{lemma}

\begin{proof}
\leanok
Follows from $T(a) \cdot T(b) = T(a+b)$ with $a = -1$, $b = 1$.
\end{proof}

\subsection{Multiplication by $x$ shifts coefficients}

\begin{lemma}[Multiplication by $x$ and $x^{-1}$ shifts coefficients]
\label{lem.fps.laure.xa}
\lean{AlgebraicCombinatorics.FPS.Laurent.T_mul_coeff, AlgebraicCombinatorics.FPS.Laurent.T_neg_one_mul_coeff}
\leantarget
\leanok
Let $\mathbf{a} = \left(a_n\right)_{n \in \mathbb{Z}}$ be a Laurent
polynomial in $K\left[\xpm\right]$. Then,
\[
  x \cdot \mathbf{a} = \left(a_{n-1}\right)_{n \in \mathbb{Z}}
  \qquad \text{and} \qquad
  x^{-1} \cdot \mathbf{a} = \left(a_{n+1}\right)_{n \in \mathbb{Z}}.
\]
\end{lemma}

\begin{proof}
From $x = \left(\delta_{i,1}\right)_{i \in \mathbb{Z}}$ and
$\mathbf{a} = \left(a_i\right)_{i \in \mathbb{Z}}$, we obtain
\[
  x \cdot \mathbf{a}
  = \left(\delta_{i,1}\right)_{i \in \mathbb{Z}}
    \cdot \left(a_i\right)_{i \in \mathbb{Z}}
  = \left(c_n\right)_{n \in \mathbb{Z}},
  \quad \text{where} \quad
  c_n = \sum_{i \in \mathbb{Z}} \delta_{i,1}\,a_{n-i}
\]
(by Definition~\ref{def.fps.laure.laupol}). For each
$n \in \mathbb{Z}$,
\begin{align*}
  c_n
  &= \sum_{i \in \mathbb{Z}} \delta_{i,1}\,a_{n-i}
   = \underbrace{\delta_{1,1}}_{= 1} a_{n-1}
     + \sum_{\substack{i \in \mathbb{Z} \\ i \neq 1}}
       \underbrace{\delta_{i,1}}_{= 0\text{ (since } i \neq 1\text{)}}
       a_{n-i} \\
  &= a_{n-1}.
\end{align*}
Thus, $x \cdot \mathbf{a} = \left(a_{n-1}\right)_{n \in \mathbb{Z}}$.

It remains to prove that
$x^{-1} \cdot \mathbf{a} = \left(a_{n+1}\right)_{n \in \mathbb{Z}}$.
Let $\mathbf{b} = \left(a_{n+1}\right)_{n \in \mathbb{Z}}$; this is
again a Laurent polynomial in $K\left[\xpm\right]$.
Applying the result just proved to
$\mathbf{b}$ (with $a_{n+1}$ in place of $a_n$) yields
\[
  x \cdot \mathbf{b}
  = \left(a_{(n-1)+1}\right)_{n \in \mathbb{Z}}
  = \left(a_n\right)_{n \in \mathbb{Z}}
  = \mathbf{a}.
\]
Dividing by the invertible element $x$, we get
$\mathbf{b} = x^{-1} \cdot \mathbf{a}$, i.e.\
$x^{-1} \cdot \mathbf{a} = \left(a_{n+1}\right)_{n \in \mathbb{Z}}$.
\end{proof}

\begin{lemma}[Multiplication on the right by $T$ shifts coefficients]
\label{lem.fps.laure.xa.right}
\lean{AlgebraicCombinatorics.FPS.Laurent.mul_T_coeff, AlgebraicCombinatorics.FPS.Laurent.mul_T_neg_one_coeff}
\leanhelper
\leanok
For any Laurent polynomial $\mathbf{a}$,
\[
  \mathbf{a} \cdot x = \left(a_{n-1}\right)_{n \in \mathbb{Z}}
  \qquad \text{and} \qquad
  \mathbf{a} \cdot x^{-1} = \left(a_{n+1}\right)_{n \in \mathbb{Z}}.
\]
\end{lemma}

\begin{proof}
Follows from Lemma~\ref{lem.fps.laure.xa} by commutativity of
multiplication.
\end{proof}

\subsection{Powers of $x$}

\begin{lemma}[$T(m) \cdot T(n) = T(m + n)$]
\label{lem.fps.laure.T-mul-T}
\lean{AlgebraicCombinatorics.FPS.Laurent.T_mul_T}
\leanhelper
\leanok
For all $m, n \in \mathbb{Z}$, we have $T(m) \cdot T(n) = T(m + n)$
in the Laurent polynomial ring. This is the fundamental multiplicative
law for Laurent monomials.
\end{lemma}

\begin{proof}
\leanok
This is the group homomorphism property of $T$: $T(a) \cdot T(b) = T(a+b)$.
\end{proof}

\begin{lemma}[$T(1)^n = T(n)$ for natural numbers]
\label{lem.fps.laure.T-one-pow}
\lean{AlgebraicCombinatorics.FPS.Laurent.T_one_pow}
\leanhelper
\leanok
For any natural number $n$, we have $T(1)^n = T(n)$ in the Laurent
polynomial ring.
\end{lemma}

\begin{proof}
\leanok
Follows from $T(1)^n = T(1 \cdot n) = T(n)$ by the group homomorphism property.
\end{proof}

\begin{proposition}[$x^k = (\delta_{i,k})_{i \in \mathbb{Z}}$ for all $k \in \mathbb{Z}$]
\label{prop.fps.laure.xk}
\lean{AlgebraicCombinatorics.FPS.Laurent.T_coeff_eq, AlgebraicCombinatorics.FPS.Laurent.T_one_zpow}
\leantarget
\leanok
We have
\[
  x^k = \left(\delta_{i,k}\right)_{i \in \mathbb{Z}}
  \qquad \text{for each } k \in \mathbb{Z}.
\]
\end{proposition}

\begin{proof}
This is proved by two-sided induction on $k$ (see \cite{detnotes},
\S 2.15 for a detailed explanation of two-sided induction).

\textbf{Base case ($k = 0$):}
$x^0 = 1 = \left(\delta_{i,0}\right)_{i \in \mathbb{Z}}$, which is the
definition of the unity element.

\textbf{Induction step from $k$ to $k+1$:}
Assume $x^k = \left(\delta_{i,k}\right)_{i \in \mathbb{Z}}$.
By Lemma~\ref{lem.fps.laure.xa} (the
$x \cdot \mathbf{a} = \left(a_{n-1}\right)_{n \in \mathbb{Z}}$ part),
\[
  x^{k+1} = x \cdot x^k
  = x \cdot \left(\delta_{i,k}\right)_{i \in \mathbb{Z}}
  = \left(\delta_{n-1,k}\right)_{n \in \mathbb{Z}}
  = \left(\delta_{n,k+1}\right)_{n \in \mathbb{Z}},
\]
since $n - 1 = k$ if and only if $n = k + 1$.

\textbf{Induction step from $k$ to $k-1$:}
Assume $x^k = \left(\delta_{i,k}\right)_{i \in \mathbb{Z}}$.
By Lemma~\ref{lem.fps.laure.xa} (the
$x^{-1} \cdot \mathbf{a} = \left(a_{n+1}\right)_{n \in \mathbb{Z}}$ part),
\[
  x^{k-1} = x^{-1} \cdot x^k
  = x^{-1} \cdot \left(\delta_{i,k}\right)_{i \in \mathbb{Z}}
  = \left(\delta_{n+1,k}\right)_{n \in \mathbb{Z}}
  = \left(\delta_{n,k-1}\right)_{n \in \mathbb{Z}},
\]
since $n + 1 = k$ if and only if $n = k - 1$.
\end{proof}

\begin{lemma}[Coefficient of $x^k$]
\label{lem.fps.laure.xk.coeff}
\lean{AlgebraicCombinatorics.FPS.Laurent.T_one_zpow_coeff}
\leanhelper
\leanok
For all $k, i \in \mathbb{Z}$, the coefficient of $x^k$ at position
$i$ is $\delta_{i,k}$.
\end{lemma}

\begin{proof}
\leanok
Immediate from Proposition~\ref{prop.fps.laure.xk}.
\end{proof}

\subsection{Proof of Proposition~\ref{prop.fps.laure.a=sumaixi} (sketched)}

\begin{definition}[Summable family of monomials]
\label{def.fps.laure.singleFamily}
\lean{AlgebraicCombinatorics.FPS.Laurent.singleFamily}
\leanhelper
\leanok
Given a Laurent series $\mathbf{a} = (a_n)_{n \in \mathbb{Z}}$, the family
of monomials $\bigl(a_g \cdot x^g\bigr)_{g \in \operatorname{supp}(\mathbf{a})}$
forms a summable family. This
construction is the key ingredient for expressing a Laurent series
as a formal infinite sum of monomials.
\end{definition}

\begin{lemma}[Coefficient of a Laurent series as a finitary sum]
\label{lem.fps.laure.coeff-finsum-single}
\lean{AlgebraicCombinatorics.FPS.Laurent.laurentSeries_coeff_eq_finsum_single}
\leanhelper
\leanok
For any Laurent series $\mathbf{a} = (a_n)_{n \in \mathbb{Z}}$ and any
$n \in \mathbb{Z}$, the $n$-th coefficient satisfies
\[
  a_n = \sum_{i \in \operatorname{supp}(\mathbf{a})}^{\mathrm{fin}}
        [x^n](a_i \cdot x^i).
\]
\end{lemma}

\begin{proof}
Follows by rewriting the left-hand side as the sum of its
monomial components and extracting the coefficient.
\end{proof}

Proposition~\ref{prop.fps.laure.a=sumaixi} follows: any Laurent
polynomial $\mathbf{a} = (a_i)_{i \in \mathbb{Z}}$ satisfies
$\mathbf{a} = \sum_{i \in \mathbb{Z}} a_i x^i$, analogously to
Corollary~\ref{cor.fps.sumakxk} but using
Proposition~\ref{prop.fps.laure.xk} instead of
Proposition~\ref{prop.fps.xk}.

\subsection{Proof of Theorem~\ref{thm.fps.laure.lauser-ring} (sketched)}

The proof is analogous to Theorem~\ref{thm.fps.laure.laupol-ring}.
The only different piece is the proof that $K\left(\left(x\right)\right)$
is closed under multiplication.

Let $\left(a_n\right)_{n \in \mathbb{Z}}$ and
$\left(b_n\right)_{n \in \mathbb{Z}}$ be two elements of
$K\left(\left(x\right)\right)$. Their product is
\[
  \left(a_n\right)_{n \in \mathbb{Z}} \cdot \left(b_n\right)_{n \in \mathbb{Z}}
  = \left(c_n\right)_{n \in \mathbb{Z}},
  \quad \text{where} \quad
  c_n = \sum_{i \in \mathbb{Z}} a_i\,b_{n-i}.
\]
We must show that $\left(c_{-1}, c_{-2}, c_{-3}, \ldots\right)$ is
essentially finite.

Since $\left(a_n\right) \in K\left(\left(x\right)\right)$, there exists
a negative integer $p$ such that $a_i = 0$ for all $i \le p$.
Similarly, there exists a negative integer $q$ such that $b_j = 0$
for all $j \le q$. Set $r := p + q$. For any integer $n \le r$:

\begin{itemize}
\item Each $i \ge p$ satisfies $n - i \le r - p = q$, hence
  $b_{n-i} = 0$.
\item Each $i < p$ satisfies $a_i = 0$.
\end{itemize}

Therefore
\[
  c_n
  = \sum_{i \in \mathbb{Z}} a_i\,b_{n-i}
  = \sum_{\substack{i \in \mathbb{Z} \\ i < p}}
    \underbrace{a_i}_{= 0} b_{n-i}
  + \sum_{\substack{i \in \mathbb{Z} \\ i \ge p}}
    a_i \underbrace{b_{n-i}}_{= 0}
  = 0.
\]
Hence all $n \le r$ satisfy $c_n = 0$, so the sequence
$\left(c_{-1}, c_{-2}, \ldots\right)$ is essentially finite.

\subsection{Helpers for the embeddings}

\begin{lemma}[Power series coefficient preservation]
\label{lem.fps.laure.ps-embed-coeff}
\lean{AlgebraicCombinatorics.FPS.Laurent.coeff_powerSeriesToLaurent}
\leanhelper
\leanok
The embedding $K[\![x]\!] \hookrightarrow K(\!(x)\!)$ preserves
coefficients: for any power series $f$ and any $n \in \mathbb{N}$,
the $n$-th coefficient of the image equals the $n$-th coefficient of $f$.
\end{lemma}

\begin{proof}
\leanok
This is a standard result relating power series
to Laurent series: the embedding preserves coefficients.
\end{proof}

\begin{lemma}[Laurent polynomial embedding is additive]
\label{lem.fps.laure.laupol-embed-add}
\lean{AlgebraicCombinatorics.FPS.Laurent.laurentPolynomialToSeries_add}
\leanhelper
\leanok
The embedding $K[\xpm] \hookrightarrow K(\!(x)\!)$ preserves
addition: for Laurent polynomials $p$ and $q$, the image of $p + q$
equals the sum of the images.
\end{lemma}

\begin{proof}
\leanok
By extensionality on coefficients: both sides have the same coefficient
at each $n \in \mathbb{Z}$.
\end{proof}

\begin{lemma}[Laurent polynomial embedding preserves 1]
\label{lem.fps.laure.laupol-embed-one}
\lean{AlgebraicCombinatorics.FPS.Laurent.laurentPolynomialToSeries_one}
\leanhelper
\leanok
The embedding $K[\xpm] \hookrightarrow K(\!(x)\!)$ sends $1$ to $1$.
\end{lemma}

\begin{proof}
\leanok
By extensionality on coefficients.
\end{proof}

\begin{lemma}[Laurent polynomial embedding preserves coefficients]
\label{lem.fps.laure.laupol-embed-coeff}
\lean{AlgebraicCombinatorics.FPS.Laurent.laurentPolynomialToSeries_coeff}
\leanhelper
\leanok
For any Laurent polynomial $p$ and any $n \in \mathbb{Z}$,
the $n$-th coefficient of the image equals $p(n)$.
\end{lemma}

\begin{proof}
\leanok
By definition of the embedding.
\end{proof}

\subsection{Helpers for Theorem~\ref{thm.fps.laure.balanced-tern-rep-uniq}}

\begin{lemma}[Existence of balanced ternary representation]
\label{lem.fps.laure.bt-exists}
\lean{AlgebraicCombinatorics.FPS.Laurent.balancedTernary_exists}
\leanhelper
\leanok
For any integer $n$, there exists a balanced ternary representation
of $n$. The proof finds $k$ large enough that $|n| \le M_k$, then
uses the $k$-bounded existence result
(Lemma~\ref{lem.fps.laure.kbounded-unique}).
\end{lemma}

\begin{proof}
\leanok
Find $k$ with $|n| \le M_k$ (using $|n| \le M_{|n|}$).
Then $n \in [-M_k, M_k]$, so by Lemma~\ref{lem.fps.laure.kbounded-unique}
there exists a $k$-bounded representation, which is converted to a
finitely supported function.
\end{proof}

\begin{lemma}[Uniqueness of balanced ternary representation]
\label{lem.fps.laure.bt-unique}
\lean{AlgebraicCombinatorics.FPS.Laurent.balancedTernary_unique}
\leanhelper
\leanok
If two balanced ternary representations have the same value $n$,
then their digit functions are equal. This follows from
Lemma~\ref{lem.fps.laure.sum-powers-three-zero} applied to the
difference of the two digit sequences.
\end{lemma}

\begin{proof}
\leanok
Apply the uniqueness result for balanced ternary digits to the two digit functions,
using that both have values in $\{-1, 0, 1\}$ and the same weighted sum.
\end{proof}

\begin{lemma}[$k$-bounded values lie in the target interval]
\label{lem.fps.laure.bt-value-mem-Icc}
\lean{AlgebraicCombinatorics.FPS.Laurent.kBoundedBTValue_mem_Icc}
\leanhelper
\leanok
For any $k$-bounded balanced ternary representation
$f \in \{-1, 0, 1\}^{k+1}$, the value
$\sum_{i=0}^{k} f(i) \cdot 3^i$ lies in the interval
$[-M_k, M_k]$ where $M_k = \sum_{i=0}^{k} 3^i$.
\end{lemma}

\begin{proof}
\leanok
Each digit $f(i) \in \{-1, 0, 1\}$ contributes at most $\pm 3^i$
to the sum. Summing gives the bounds
$-\sum 3^i \le \text{value} \le \sum 3^i$.
\end{proof}

\begin{lemma}[Balanced ternary injectivity]
\label{lem.fps.laure.bt-injective}
\lean{AlgebraicCombinatorics.FPS.Laurent.kBoundedBTValue_injective_on}
\leanhelper
\leanok
The evaluation map $f \mapsto \sum_{i=0}^{k} f(i) \cdot 3^i$ is
injective on the set of $k$-bounded balanced ternary representations.
\end{lemma}

\begin{proof}
By strong induction on $k$: the lowest digit is determined by the value
modulo $3$, and the remaining digits are determined by induction on
the quotient.
\end{proof}

\begin{lemma}[Digit differences are bounded]
\label{lem.fps.laure.bt-diff-abs}
\lean{AlgebraicCombinatorics.FPS.Laurent.btDigits_diff_abs_ge_one, AlgebraicCombinatorics.FPS.Laurent.btDigits_diff_abs_le_two}
\leanhelper
\leanok
If $a, b \in \{-1, 0, 1\}$, then $|a - b| \le 2$.
If additionally $a \ne b$, then $|a - b| \ge 1$.
\end{lemma}

\begin{proof}
\leanok
By case analysis on $a$ and $b$.
\end{proof}

\begin{lemma}[Equivalence of balanced ternary representation types]
\label{lem.fps.laure.bt-equiv}
\lean{AlgebraicCombinatorics.FPS.Laurent.balancedTernary_equiv}
\leanhelper
\leanok
There is a type equivalence between the finitely-supported
balanced ternary representation (using $\mathbb{N} \to \mathbb{Z}$
with digits in $\{-1,0,1\}$ and finite support) and the inductive-style
balanced ternary representation. The conversions are inverse
to each other.
\end{lemma}

\begin{proof}
The forward direction maps each digit $d$ to its balanced ternary digit,
the inverse maps each balanced ternary digit to its integer value.
Round-trip identities follow by checking both compositions are the identity.
\end{proof}
